\par\Needspace{6\baselineskip}\medskip
\SourceLeaf{142}\pdfbookmark[2]{102.}{article-102}
\begin{center}102.\end{center}
Since $y$ must be positive in the hyperbola just as in the ellipse, the solution of equation 16 is unambiguous here as well\footnote{It scarcely needs mentioning that our Table II can be used to solve this equation for the hyperbola just as for the ellipse, as long as $h$ remains within its limits.}. Equation $16^*$, however, requires a slightly different argument here than for the ellipse. The theory of equations readily shows that for a positive value of $H$\HistoricalNoteDoubleUnbracketed{The quantity $H$ can evidently become negative only if $\zeta>\tfrac16$. Such a value of $\zeta$ would correspond to a value of $z$ greater than $2{,}684$, far beyond the limits of this method.} this equation, if it has any positive real root at all, has one negative root and two positive roots. The two positive roots are either equal, with both $=\tfrac16\BareRoot5-\tfrac16=0{,}20601$, or one is greater and the other smaller than this limit. We now show that in our problem the larger root must always be taken. The problem rests on the assumption that $z$ is small and, at any rate, no greater than $0{,}3$, so that the third table can still be used. If in equation $13^*$ we replace $M$ with $Y\BareRoot{(L+z)}$, we obtain
\[
Y+1=(L+z)Z>(1+z)Z,
\]
\noindent or
\[
Y>\tfrac13-\frac4{3\mathpunct{.}5}z+\frac{4\mathpunct{.}6}{3\mathpunct{.}5\mathpunct{.}7}zz-\frac{4\mathpunct{.}6\mathpunct{.}8}{3\mathpunct{.}5\mathpunct{.}7\mathpunct{.}9}z^3+\text{etc.},
\]
\noindent from which it readily follows that for values of $z$ as small as we assume here, we must always have $Y>0{,}20601$. Indeed, we find by calculation that equality of $(1+z)Z$ with this limit would require $z=0{,}79858$. Such large values of $z$ lie well beyond the intended scope of our method.
